\documentclass[11pt]{article}
\usepackage[utf8]{inputenc}
\usepackage{amsmath, amssymb, amsthm}
\usepackage{geometry}
\geometry{margin=1in}

\newtheorem{theorem}{Theorem}
\newtheorem{definition}{Definition}
\newtheorem{proposition}{Proposition}

\title{Probability Density Functions}
\author{math-foundations / concept 24-pdf}
\date{}

\begin{document}
\maketitle

\subsection*{First, an example}
Before the abstract definition, fix the standard normal PDF
$\phi(x) = \frac{1}{\sqrt{2\pi}}\,e^{-x^2/2}$ in mind.
At $x = 0$ its height is
$\phi(0) = 1/\sqrt{2\pi} \approx 0.399$.
That value is greater than nothing in particular --- but observe: a PDF
\emph{can} exceed $1$ at a point. What must equal $1$ is the
\emph{area} under the curve, not the heights:
\[
    \int_{-\infty}^{\infty} \phi(x)\,dx \;=\; 1.
\]
\emph{Read aloud:} the integral from minus infinity to infinity of phi of x, d x, equals one.
Density measures probability per unit length; the bookkeeping rule is
on the integral, not on pointwise values.

\section{Definitions}

\begin{definition}[Probability density function]
Let $X$ be a real-valued random variable on a probability space $(\Omega, \mathcal{F}, P)$. A non-negative, Lebesgue-integrable function $f_X : \mathbb{R} \to [0, \infty)$ is called the \emph{probability density function} (PDF) of $X$ if for every Borel set $B \subseteq \mathbb{R}$,
\[
    P(X \in B) \;=\; \int_B f_X(x)\,dx.
\]
\emph{Read aloud:} the probability that X lies in B equals the integral over B of f sub X of x, d x.
Equivalently, the law of $X$ is absolutely continuous with respect to Lebesgue measure, and $f_X = dP_X/d\lambda$ is the Radon--Nikodym derivative.
\end{definition}

\begin{definition}[CDF--PDF relation]
The cumulative distribution function $F_X(x) = P(X \le x)$ satisfies
\[
    F_X(x) = \int_{-\infty}^{x} f_X(t)\,dt, \qquad f_X(x) = F_X'(x) \text{ a.e.}
\]
\emph{Read aloud:} F sub X of x equals the integral from minus infinity to x of f sub X of t, d t; and f sub X of x equals F sub X prime of x almost everywhere.
\end{definition}

\begin{proposition}[Normalization]
Any PDF $f_X$ satisfies $\int_{-\infty}^{\infty} f_X(x)\,dx = 1$.
\end{proposition}

\begin{proof}
$\int_\mathbb{R} f_X = P(X \in \mathbb{R}) = 1$.
\end{proof}

\section{The Gaussian PDF}

Define the \emph{standard Gaussian density}
\[
    \phi(x) \;=\; \frac{1}{\sqrt{2\pi}} \exp\!\left(-\frac{x^2}{2}\right).
\]
\emph{Read aloud:} phi of x equals one over square-root of two pi, times e to the minus x squared over two.
Clearly $\phi > 0$, so to show $\phi$ is a PDF we need only verify normalization.

\begin{theorem}[Gaussian normalization]
\label{thm:gauss}
The Gaussian PDF $\phi$ integrates to $1$ over $\mathbb{R}$:
\[
    \int_{-\infty}^{\infty} \phi(x)\,dx \;=\; 1.
\]
\emph{Read aloud:} the integral from minus infinity to infinity of phi of x, d x, equals one.
\end{theorem}

\begin{proof}
Let $I := \int_{-\infty}^{\infty} e^{-x^2/2}\,dx$. The integral is finite by comparison with a step function (or by direct integrability check). We compute $I^2$ via Fubini's theorem on the product:
\[
    I^2 \;=\; \left(\int_{-\infty}^{\infty} e^{-x^2/2}\,dx\right)\!\left(\int_{-\infty}^{\infty} e^{-y^2/2}\,dy\right)
    \;=\; \iint_{\mathbb{R}^2} e^{-(x^2 + y^2)/2}\,dx\,dy.
\]
\emph{Read aloud:} I squared equals the product of two one-dimensional Gaussian integrals, which equals the double integral over R-two of e to the minus x squared plus y squared over two, d x d y.
Switch to polar coordinates $x = r\cos\theta$, $y = r\sin\theta$, with Jacobian $r$:
\[
    I^2 \;=\; \int_{0}^{2\pi}\!\int_{0}^{\infty} e^{-r^2/2}\,r\,dr\,d\theta.
\]
\emph{Read aloud:} I squared equals the integral from zero to two pi, of the integral from zero to infinity of e to the minus r squared over two, times r, d r, d theta.
The inner integral evaluates by substitution $u = r^2/2$, $du = r\,dr$:
\[
    \int_{0}^{\infty} r\,e^{-r^2/2}\,dr \;=\; \int_{0}^{\infty} e^{-u}\,du \;=\; 1.
\]
\emph{Read aloud:} the integral from zero to infinity of r e to the minus r squared over two, d r, equals the integral from zero to infinity of e to the minus u, d u, which equals one.
Therefore
\[
    I^2 \;=\; \int_{0}^{2\pi} 1\,d\theta \;=\; 2\pi,
\]
\emph{Read aloud:} I squared equals the integral from zero to two pi of one, d theta, which equals two pi.
so $I = \sqrt{2\pi}$ (taking the positive root, as $e^{-x^2/2} > 0$). Hence
\[
    \int_{-\infty}^{\infty} \phi(x)\,dx \;=\; \frac{1}{\sqrt{2\pi}} \cdot \sqrt{2\pi} \;=\; 1. \qedhere
\]
\emph{Read aloud:} the integral of phi of x, d x, equals one over square-root two pi, times square-root two pi, which equals one.
\end{proof}

\begin{proposition}[Gaussian CDF]
The standard Gaussian CDF $\Phi(x) = \int_{-\infty}^{x} \phi(t)\,dt$ admits no elementary closed form, but
\[
    \Phi(x) \;=\; \tfrac{1}{2}\Bigl(1 + \operatorname{erf}\!\bigl(x/\sqrt{2}\bigr)\Bigr),
\]
\emph{Read aloud:} capital Phi of x equals one half times the quantity one plus the error function of x over square-root two.
where $\operatorname{erf}(z) = \frac{2}{\sqrt{\pi}}\int_0^z e^{-t^2}\,dt$.
\end{proposition}

\section{Remark on densities versus probabilities}

A PDF can exceed $1$. For example, the uniform PDF on $[0, 1/100]$ is $f(x) = 100$ on its support, yet $\int f = 1$. Density is probability \emph{per unit length}, not probability itself; pointwise probability $P(X = x) = 0$ for any continuous random variable.

\end{document}
